\documentclass[10pt,a4paper]{article}
\usepackage[margin=15mm,top=14mm,bottom=16mm]{geometry}
\usepackage[T1]{fontenc}
\usepackage[utf8]{inputenc}
\usepackage{lmodern,microtype,amsmath,amssymb,graphicx,booktabs,array,xcolor}
\usepackage[font=small,labelfont=bf,skip=4pt]{caption}
\usepackage[colorlinks=true,urlcolor=blue!50!black,linkcolor=blue!50!black]{hyperref}
\hypersetup{pdftitle={Singlet coverings in two and three dimensions},
  pdfsubject={Square--octagon RVB, octochlore coverings, and a possible layer construction}}
\usepackage{fancyhdr}
\pagestyle{fancy}
\fancyhf{}
\renewcommand{\headrulewidth}{0pt}
\fancyfoot[L]{\footnotesize Singlet coverings: collaborator notes \quad 22 October 2026}
\fancyfoot[R]{\footnotesize\thepage/3}
\setlength{\parindent}{0pt}
\setlength{\parskip}{5pt}
\setlength{\abovedisplayskip}{5pt}
\setlength{\belowdisplayskip}{5pt}
\setlength{\textfloatsep}{5pt}
\newcommand{\sect}[1]{\vspace{5pt}{\large\bfseries\color{blue!45!black}#1}\par\vspace{2pt}}
\newcommand{\ket}[1]{\lvert#1\rangle}
\newcommand{\eps}{\epsilon}
\newcommand{\xor}{\mathbin{\oplus}}
\begin{document}

{\LARGE\bfseries Singlet coverings in two and three dimensions}\par
{\normalsize Square--octagon RVB, octochlore coverings, and a possible layer construction}

\sect{1. Two-dimensional model and the two dummy-link arrangements}
The underlying graph is the trivalent square--octagon lattice. Its gray squares $P$ are
drawn upright; a physical spin $1/2$ sits at the midpoint of each square edge
(blue). Each dummy link carries \emph{two} distinct endpoint spins (orange).
The three slots around a graph vertex form a triangle, which contains zero or
one singlet. Every physical spin is covered exactly once globally.

\begin{center}
\includegraphics[width=.93\linewidth]{upright-geometry.pdf}
\end{center}
\vspace{-6pt}
\captionof{figure}{Gray lines define the base graph; green triangle edges are allowed
singlet bonds. The gray squares have blue spins on their edges. A cell contains
four blue spins and four orange endpoint spins.}

On a dummy link either its two endpoint spins form an \textbf{internal singlet},
or \textbf{both endpoints pair into their own triangles}. There is no singly
occupied case and no free spin. For the internal case each neighbouring triangle
may be empty or contain its blue--blue singlet: four possibilities. For the
external case each endpoint chooses one of two blue spins: another four.

\begin{center}
\includegraphics[width=\linewidth]{dummy-covering-table.pdf}
\end{center}
\vspace{-6pt}
\captionof{table}{All eight local completions and their square-vertex loop patterns.
Green segments are singlets; dotted green continuations pair a blue spin in its
other triangle. After contraction the ports are $a=N$, $b=W$, $c=E$, $d=S$.
The displayed bits are relative to the representative local reference $0000$,
with the orange spins internally paired. This reference extends to a global covering.}

Writing $a,b,c,d=1$ when an ordinary slot is used in its displayed triangle, and
$t=1$ when a dummy endpoint is used there, the local constraints are
\[
 a\xor b=t=c\xor d,
 \qquad a\xor b\xor c\xor d=0.
\]
In this representative reference these are also the loop bits. For any other
reference, first take the difference of ownership bits modulo two.

The conventional RVB state sums the normalized singlet products with equal coefficients:
\[
 \ket{\Psi_{2D}}=\mathcal N^{-1/2}\sum_{D\ {
m allowed}}\ket{D},
 \qquad
 \ket{s_{ij}}=\frac{\ket{\uparrow_i\downarrow_j}-\ket{\downarrow_i\uparrow_j}}{\sqrt2}.
\]
Singlet arrows run counterclockwise around each triangle; dummy arrows follow
the positive simulation axes. This fixes relative phases. No additional loop-number
sign is included. Equal coefficients do not make overlapping singlet coverings orthogonal.

\newpage
\sect{2. Coverings $\longrightarrow$ loops $\longrightarrow$ the square lattice}
Choose one valid reference covering. An ordinary-edge spin has exactly one owner;
the two dummy endpoints are either both used by triangles or neither is. Ownership
differences agree at both ends of every edge and have even parity at each base vertex.
Since the base graph has degree three, the selected edges form closed loops [1].

\begin{center}
\includegraphics[width=\linewidth]{corrected-mapping-upright.pdf}
\end{center}
\vspace{-6pt}
\captionof{figure}{One checked covering and its reference-relative loops. The relevant
gray squares are upright in \emph{all three panels}. The reference has internal dummy
singlets and alternating occupied/empty triangles around each square. Boundary singlet
stubs continue periodically. Dummy edges shrink to vertices; square edges survive.}

Here ``cutting/erasing'' a dummy edge means \textbf{contracting its two graph endpoints},
not deleting its physical spins from the RVB state. The result is the regular square
graph. Writing $n$ for reference-difference bits, the erased bit is recovered as
$n_{\rm dummy}=n_a\xor n_b=n_c\xor n_d$, so the map is a
bijection of even-edge configurations, including winding sectors. At a four-leg
vertex retain the inherited $a$--$b$, $c$--$d$ routing shown in Table~1.
These reference loops label coverings; bra--ket transition loops instead determine
singlet overlaps. Consequently the map alone does not determine quantum correlations.

\sect{3. Bond correlations: numerical evidence and diagnosis}
We measure the connected physical bond covariance
$C_B(r)=\langle B_0B_r\rangle-\langle B_0\rangle\langle B_r\rangle$,
$B=\mathbf S_i\!\cdot\!\mathbf S_j$, for a square-edge-spin bond inside a triangle.
It equals the connected covariance of the singlet projector $1/4-B$.

\begin{center}
\includegraphics[width=\linewidth]{correlations-2d.pdf}
\end{center}
\vspace{-6pt}
\captionof{figure}{Our model's computed bond correlations, not published curves.
An exact $D=3$ PEPS is contracted with infinite-plane CTMRG ($\chi=64$), comparing
two truncation cutoffs. Shading at left marks the fit window $r=12$--18.}

Following PEPS boundary and sector diagnostics [2,3], we compare cutoffs, fit windows,
and exact finite-cylinder checks. Within each fixed parity sector, transfer modes
coupled to $B$ give $\xi_a=-1/\log|\lambda_a/\lambda_0|$.
Fits to $Ae^{-r/\xi}$ and $Ar^{-p}e^{-r/\xi}$ give
$\xi=0.82$ and $0.86$ cells. The local inverse length
$\kappa_{\rm eff}=\log|C_B(r)/C_B(r+1)|$ tends to $1/\xi>0$ for an exponential
and to zero for a pure power law. The resolved data support a short exponential-like
scale, \textbf{$\xi\sim0.8$--$0.9$ simulation-cell translations}; an 8\% cutoff
change at $r=18$ prevents a precise bulk claim. One translation is a diagonal
$(1,-1)$ step on the upright contracted grid. Orthogonal-covering occupation
correlations instead have exactly finite range; they are a different observable and state.

{\footnotesize
[1] M.~Iqbal, D.~Poilblanc and N.~Schuch, \emph{Semionic resonating valence bond states}
(2014), \href{https://arxiv.org/abs/1407.7773}{arXiv:1407.7773}. Reference-arrow construction;
our state has no semionic sign.\par
[2] N.~Schuch et al., \emph{Resonating valence bond states in the PEPS formalism}
(2012), \href{https://arxiv.org/abs/1203.4816}{arXiv:1203.4816}.\par
[3] D.~Poilblanc et al., \emph{Topological and entanglement properties of resonating
valence bond wave functions} (2012), \href{https://arxiv.org/abs/1202.0947}{arXiv:1202.0947}.
The papers motivate the diagnostics; they do not calculate this square--octagon model.}

\newpage
\sect{4. Three-dimensional octochlore covering model}
Octahedron centres form a simple-cubic lattice, $v\in\mathbb Z^3$, with physical
sites at the shared corners $v\pm\hat e_\mu/2$. Centres are not spin sites [4].
A diameter dimer joins opposite vertices; a trimer occupies one of eight triangular
faces. Every site has exactly one owner. The local catalogue from the supplied notes is
\begin{center}
\begin{tabular}{@{}p{.54\linewidth}cl@{}}
\toprule
Allowed local objects & Number & Incoming support \\
\midrule
Any subset of $D_x,D_y,D_z$ & $1+3+3+1=8$ & Unions of opposite pairs \\
Exactly one face trimer & $8$ & One site from each axis pair \\
\bottomrule
\end{tabular}
\end{center}
Two complementary trimers are excluded: otherwise the all-six-incoming support
would cease to specify a unique covering. Diameter chords may cross at a centre,
which carries no spin. With ownership bits $a^{\pm}_{v,\mu}$, the sixteen supports obey
\[
 a^+_{v,x}\xor a^-_{v,x}
 =a^+_{v,y}\xor a^-_{v,y}
 =a^+_{v,z}\xor a^-_{v,z}.
\]

For the quantum extension put \textbf{spin 1} at every site. In the Cartesian basis,
the normalized two- and three-spin singlets are
\[
 d_{\alpha\beta}=\delta_{\alpha\beta}/\sqrt3,
 \qquad t_{\alpha\beta\gamma}=\eps_{\alpha\beta\gamma}/\sqrt6,
 \qquad \ket{\Psi_{3D}}\propto\sum_{C\ {
m allowed}}\ket{C}.
\]
Order every trimer by its $x,y,z$ vertices, independent of face signs; this fixes
its antisymmetric phase. This is a specified trial state, with no assumed parent
Hamiltonian. Current finite-box spin calculations reach $6^3$ centres (648 spins),
but do not yet establish a bulk correlation length.

\sect{5. Layer geometry and a route to investigate foliation}
\begin{center}
\includegraphics[width=\linewidth]{layer-relation.pdf}
\end{center}
\vspace{-6pt}
\captionof{figure}{Equatorial perimeter edges of one $z$ layer form a square graph.
Splitting each shared corner into two base vertices joined by a dummy edge gives
square--octagon. Contracting the dummy edges reverses the construction.}

This is an exact \textbf{geometric resolution of the equatorial layer}, rather than
a literal projection of all octahedron edges: the apices and their spokes have been
omitted. The equatorial square graph has physical spins at its vertices. The split
vertices in the last panel are graph vertices; the triangle-decorated 2D model places
its physical spins on the associated slots. This operation alone is not a spin-space map.

There is also a useful \textbf{planar constraint relation}. For differences $n_e$
from a fixed 3D reference covering, define
$p_\mu=n_{+\mu}\xor n_{-\mu}$. The cage rule $p_x=p_y=p_z$ gives an even edge
pattern in every coordinate plane of the \emph{cubic-centre graph}. Vertical parity
couples these planes: an $xy$-plane turn has $p_x=p_y=1$ and requires odd vertical
incidence. Two matching plaquette loops on adjacent planes, joined by four vertical
corner edges, form an allowed cube cage. A lone planar plaquette does not.

These facts motivate a layer construction, but \textbf{quantum foliation remains a
proposal}. It requires an explicit state map and interlayer coupling, matching
amplitudes and singlet phases. In particular, the 3D spin-1 space and 2D spin-$1/2$
slots differ: representing spin 1 by a symmetric pair does not accommodate the
antisymmetric internal dummy singlet of the 2D model.

{\footnotesize
Source of the sixteen-state rule: J.~Garre Rubio, \emph{Put the objects on octahedra /
Octahedral coverings and X-cube cage nets}, supplied notes, 19 October 2026.\par
[4] O.~Benton and R.~Moessner, \emph{Topological Route to New and Unusual Coulomb Spin
Liquids}, Phys.~Rev.~Lett.~127, 107202 (2021),
\href{https://arxiv.org/abs/2103.10817}{arXiv:2103.10817}. Geometry reference;
the sixteen-state rule and the singlet extension are specified above.}

\end{document}
